\section{Introduction}

Psychiatric diagnoses group biologically heterogeneous patients under a single label, and two
patients meeting criteria for the same disorder may share little in their cognition, biology or
course. Dimensional frameworks were introduced to replace these categories with continuous axes
that cut across diagnoses~\citep{insel2010rdoc,kotov2017hitop,ruggero2019hitop}, and a general
factor of psychopathology---the tendency for symptoms to co-occur and severity to load
together---has become their organizing backbone~\citep{caspi2014pfactor,lahey2017general}. A
dimensional alternative, however, is only as good as the \emph{measurement} that produces it: to
be credible across cohorts it must be comparable between sites, faithful to the mixed data types
of a clinical assessment, honest about the missingness that pervades real registries, and stable
enough that an axis defined once keeps its meaning when it is later validated.

Where biology has entered these maps, it has almost always entered as a \emph{correlate}---a
biomarker regressed on a symptom dimension, or a stratification computed after the fact---rather
than as a coordinate of the space itself~\citep{marquand2016,wolfers2018,segal2023}. This choice
has a consequence that is easy to miss. If biology is only ever compared to severity, one cannot
ask whether a biological axis is \emph{separable} from severity, because the machinery that would
answer the question has been assumed away. The immunometabolic burden of severe mental
illness---weight, adiposity, dysglycaemia, inflammation---is a natural test case: it is elevated
transdiagnostically, tracks medication and chronicity, and has been argued to form a partly
distinct dimension of illness~\citep{vancampfort2015mets,pillinger2017,perry2019,penninx2024imd,milaneschi2020,lamers2020}.
Yet whether that dimension stands \emph{apart from} the general severity axis, rather than
riding along with it, has been described qualitatively but never measured inside a single
transdiagnostic map.

Here we build such a map. We fit one global, missingness-aware Bayesian bifactor model to the
harmonized baseline assessment of $N=9{,}013$ patients with bipolar disorder, schizophrenia or
major depression from three national cohorts, with diagnosis withheld from model fitting
throughout, and place routinely-collected biology \emph{inside} the factor space as a co-equal
latent axis alongside cognition, sleep and the symptom dimensions. We then interrogate the fixed
map with a single question asked of every validation stage: \emph{what stands apart, what stays,
and what predicts?} The answer converges on one axis. Immunometabolic load is nearly independent
of general burden when the model is free to entangle them, is the most temporally durable
dimension of the eight, and marks the worst-prognosis pole of the patient space---an
``immunometabolic island'' that a diagnosis or a severity score does not resolve.

We are explicit about what is and is not new. That immunometabolic biology matters in severe
mental illness, and forms a partly separable dimension, is established; our finding
\emph{converges} with that independent literature, which is why we report it as relative, measured
independence rather than a new biological claim. What is new is the map: to our knowledge the
first diagnosis-blind dimensional atlas spanning these three disorders that carries biology as a
co-equal axis, and the measurement discipline---no imputation, mixed likelihoods, uncertainty
propagated end-to-end, the map frozen before validation---that turns ``biology is separable from
severity'' from a modelling assumption into an estimated quantity.
